\documentclass[11pt,a4paper]{article}

% Packages
\usepackage[utf8]{inputenc}
\usepackage[T1]{fontenc}
\usepackage[french]{babel}
\usepackage{geometry}
\usepackage{graphicx}
\usepackage{float}
\usepackage{titlesec}
\usepackage{enumitem}
\usepackage{xcolor}
\usepackage{hyperref}
\usepackage{fancyhdr}
\usepackage{tcolorbox}
\usepackage{amsmath}

% Page setup
\geometry{margin=2.5cm}
\pagestyle{fancy}
\fancyhf{}
\rhead{\thepage}
\lhead{Technical Portfolio}
\renewcommand{\headrulewidth}{0.4pt}

% Colors
\definecolor{primarycolor}{RGB}{0,51,102}
\definecolor{accentcolor}{RGB}{204,102,0}

% Section formatting
\titleformat{\section}
{\color{primarycolor}\Large\bfseries}
{\thesection}{1em}{}[\titlerule]

\titleformat{\subsection}
{\color{accentcolor}\large\bfseries}
{\thesubsection}{1em}{}

% Hyperref setup
\hypersetup{
    colorlinks=true,
    linkcolor=primarycolor,
    urlcolor=primarycolor,
    citecolor=primarycolor
}

\begin{document}

% ===============================
% Title Page
% ===============================
\begin{titlepage}
    \centering

    % Top spacing
    \vspace*{2.5cm}

    % Main title
    {\Huge\bfseries\color{primarycolor}
    Technical Portfolio\par}

    \vspace{0.6cm}

    % Horizontal rule
    \rule{0.6\textwidth}{0.8pt}

    \vspace{0.8cm}

    % Subtitle
    {\Large\bfseries
    Internship Projects and Personal Studies\par}

    \vspace{1.8cm}

    % Description box
    \begin{tcolorbox}[
        width=0.85\textwidth,
        colback=gray!8,
        colframe=primarycolor,
        boxrule=0.6pt,
        arc=2mm
    ]
    \centering
    This portfolio presents a selection of academic, industrial, and personal projects,
    carried out as part of internships and self-directed studies, with a particular focus on
    numerical simulation, thermomechanical analysis, and mechanical design.
    \end{tcolorbox}

    \vfill

    % Bottom information
    {\large
    \textbf{Ali Abouelazz} \\[0.3cm]
    

\end{titlepage}


\tableofcontents
\newpage

% ============================================
\section{General Electric Vernova}
% ============================================

\subsection{Project 1: Stiffness Measurement Station}

\subsubsection{Objective}

Design (mechanical and software components) of a stiffness measurement station for the finger contact springs of the new series of live tank circuit breakers. The purpose of this station is to ensure the conformity of the finger contact springs after assembly in the housing.

\subsubsection{Execution}

My responsibility was to develop the station in its entirety, both mechanically and in terms of software. I began by defining requirements with the various stakeholders (quality department, operators, design office). These needs were translated into system-level requirements, then broken down into subsystem-level requirements and further into component-level requirements (top-down design approach).

Once the component-level requirements were defined, I began the design phase while simultaneously working with my supervisor to identify the simulations needed to validate the measurement accuracy requirements. After dimensional validation, I produced the engineering drawings, which were then sent to the procurement department.

The software component was developed in parallel. I created a graphical user interface for data visualization and control of the force sensor, then synchronized it with the site database for data logging.

\begin{figure}[H]
    \centering
    \includegraphics[width=1\textwidth]{media/image9.png}
    \caption{Mechanical design of the measurement station}
    \label{fig:ge1}
\end{figure}

\begin{figure}[H]
    \centering
    \includegraphics[width=1\textwidth]{media/image10.png}
    \caption{Software interface and data acquisition system}
    \label{fig:ge2}
\end{figure}



\subsection{Project 2: Power Transmission Shaft Sizing}

\subsubsection{Objective}

Sizing of a redesigned power transmission shaft for the opening mechanism of GCB (Ground Circuit Breaker) circuit breakers. The previous system was costly and sensitive to coaxiality defects.

\subsubsection{Execution}

I sized a new shaft design by first simulating the various proposals made by a colleague. I iterated over different dimensions and made minor modifications to achieve the required safety factor without exceeding it. The redesign aimed to simplify machining operations and, most importantly, to reduce the amount of material used, which represented a significant cost for the company.

\begin{figure}[H]
    \centering
    \includegraphics[width=1\textwidth]{media/image99.png}
    \caption{FEA analysis of the initial design}
    \label{fig:ge5}
\end{figure}

\subsection{Project 3: Stress Analysis of SF6 Gas Chamber}

\subsubsection{Objective}

Determine the mechanical stress in the SF6 gas chamber due to the gas pressure, which was increased in order to extend the service life of the tulip contacts. The objective is to ensure that the safety factor remains above 5 as specified in the design requirements (CDC 912625/CODAP).

\subsubsection{Execution}

The load was applied to the inner walls to simulate a static pressure of $p = 1.6$ MPa. Material used: A-S7G03.

\begin{figure}[H]
    \centering
    \includegraphics[width=1\textwidth]{media/image1.png}
    \caption{FEA model of the gas chamber}
    \label{fig:ge7}
\end{figure}

\newpage


% ============================================
\section{Caterpillar: Final Year Internship}
% ============================================

\subsection{Objective and Context}

The Caterpillar D5 bulldozer had been accumulating field failures on its chain tensioner assembly, grease leaks, cracked plugs, worn retainer plates, premature piston-seal failures. The product-support team had flagged the problem but the failure modes had never been quantified, and the dominant root causes were still unclear. My internship mission was to close that loop: diagnose the system end-to-end, validate the diagnosis with simulation, and deliver redesign concepts that addressed the root causes while remaining compatible with manufacturing and assembly constraints.

The project followed a full \textbf{DMAIC cycle}. Scope was restricted to the D5 variants (D5~LGP, D5~XL) and specifically to the tensioner group . The D4, D6, D7, John~Deere~700L and Komatsu~D51~PX systems were used as benchmark references. Deliverables were 3D CAD redesigns validated in FEA, not physical prototypes. Tools: Creo Parametric for CAD, ANSYS Mechanical for structural FEA, Python for data processing.

\begin{figure}[H]
    \centering
    % Top: one wide figure
    \includegraphics[width=0.95\textwidth]{media/cat/13.png}
    
    \vspace{0.5em}
    
    % Bottom: two square figures side by side
    \begin{minipage}[t]{0.47\textwidth}
        \centering
        \includegraphics[width=\linewidth]{media/cat/14.png}
    \end{minipage}
    \hfill
    \begin{minipage}[t]{0.47\textwidth}
        \centering
        \includegraphics[width=\linewidth]{media/cat/15.png}
    \end{minipage}
    
    \caption{D5 track tensioner: cross-section showing the main components involved in the failure modes (spring tube, piston, piston seal, retainer, cylinder, plug). Grease pressure acts on the piston base to tension the track; the plug seals the rear of the cylinder and is bolted to the housing.}
    \label{fig:cat_system}
\end{figure}

\subsection{Automated Failure-Mode Classification of Warranty Reports}

Caterpillar maintains a centralised warranty database: every time a dealer repairs a machine under warranty, they log the failed part, the machine hours, a short \textit{comment} field and a longer \textit{claim story}. Filtering by the tensioner part numbers returned records  \textbf{in the order of thousands}  a dataset large enough to make statistical analysis meaningful, but far too large to read manually. The real obstacle was that dealers worldwide write in their native language: English, German, Polish, Japanese, Spanish, sometimes with abbreviations and regional shorthand. Per-component \textit{frequency} was easy to compute from the part number alone, but extracting the actual \textit{failure mode} from free text required something smarter.

My first attempt used the Gemini API: feed each concatenated \texttt{comment + claim\_story} to the model with a system prompt asking it to pick one mode from a fixed list. It worked well, but every call sent customer-written warranty text to a third-party server. Even though the raw text did not mention "Caterpillar" or part numbers, sending it out over an API was not acceptable for this use case. Deploying a local LLM server (Ollama-style) would have solved the privacy issue but required IT approvals I was not going to get in time.

I pivoted to a lighter, fully local approach based on \textbf{semantic similarity with Sentence-BERT}. SBERT embeds a sentence into a dense high-dimensional vector whose geometry captures semantic meaning  two sentences expressing the same idea in different words (or different languages) land close together. I embedded each of the eight candidate failure modes (\textit{grease leak, broken threads, detached part, loose part, broken/cracked part, corroded part, clogged, damaged threads/grooves}) once, then embedded every claim story, and classified each one by the argmax of cosine similarity:

\[
\mathrm{FM}(t) \;=\; \arg\max_i \; \frac{\mathbf{v}_t \cdot \mathbf{v}_{\mathrm{MD}_i}}{\|\mathbf{v}_t\|\,\|\mathbf{v}_{\mathrm{MD}_i}\|}
\]

A similarity threshold was applied below which the case was marked null, so borderline or uninformative claim stories would not pollute the statistics. The \texttt{all-MiniLM-L6-v2} model runs comfortably on CPU and handled the multilingual content well enough that manual translation was unnecessary. The entire pipeline ran on my workstation, no data left the machine. I validated the output against 50 manually-labelled samples and the classifier hit \textbf{48/50 = 96\%} accuracy, which was well within acceptable bounds for prioritisation work.

Crossing the resulting failure-mode distribution with each component's warranty cost revealed a clear ranking of critical failures: leakage at the piston/seal interface, leakage at the fill-valve/cylinder interface, plug loosening and cracking, and leakage at the relief-valve interface. Equally striking was what the data \textit{did not} contain: zero reports of the intentional fuse system actuating, despite multiple plug-cracking cases. That contradiction became the focus of the next phase.

\begin{figure}[H]
    \centering
    % Top: two square figures side by side
    \begin{minipage}[t]{0.45\textwidth}
        \centering
        \includegraphics[width=\linewidth]{media/cat/19.png}
    \end{minipage}
    \hfill
    \begin{minipage}[t]{0.51\textwidth}
        \centering
        \includegraphics[width=\linewidth]{media/cat/20.png}
    \end{minipage}
    
    \vspace{0.5em}
    
    % Bottom: one wide figure
    \includegraphics[width=0.95\textwidth]{media/cat/21.png}
    
    \caption{Failure-mode distribution per component after automated classification of the  warranty records. Grease leak dominates across seals, valves and plug.}
    \label{fig:cat_failuremodes}
\end{figure}

\subsection{FEA of the Fuse System: Why Plugs Were Cracking}

The tensioner is built with an intentional mechanical fuse, a deformable steel bar backed by an O-ring  designed to open a leak path \textit{before} any other component reaches its yield limit. If the fuse was doing its job, cracked plugs should not exist in the warranty record. They did. Something in the actual stress response of the assembly was violating the intent of the design, and I needed an FEA model to find out what.

I built a coupled model of the fuse assembly in ANSYS Mechanical, applying grease pressure directly to the internal surfaces and pre-tensioning the bolts and plug per the design drawings. The first version modelled the O-ring explicitly as a hyperelastic body the theoretically correct choice. In practice the O-ring deformation became so large that mesh elements distorted past ANSYS's convergence tolerance, and the simulation crashed above 24 pressure units, far below the regime of interest (130--250 units).

Rather than fight the solver with finer meshes and remeshing hacks, I reformulated the leak criterion \textit{geometrically}. Per the Parker O-ring Handbook, an elastomeric seal is guaranteed to hold as long as its compression exceeds 5.7\%, and is guaranteed to leak below 0\% (loss of contact). If I remove the O-ring from the simulation entirely and track the local vertical displacement $d_{z,k}$ of each node along the seal contour, I can compute a per-node compression:

\[
C_{\%,k} \;=\; \frac{t_o - t_{c,k}}{t_o}, \qquad t_{c,k} = d_{z,k} + g_d
\]

where $t_o$ is the free O-ring thickness and $g_d$ is the gland depth. Aggregating over all contour nodes gives \textbf{guaranteed sealing} when $\min_k C_{\%,k} > 5.7\%$ and \textbf{certain leakage} when $\overline{C_\%} < 0\%$. With the O-ring removed, the model meshed cleanly and ran stable all the way to plastic yield, while still producing a physically meaningful leak prediction. The model also exploited the cylindrical symmetry of the assembly, halving the node count and the solve time.

The results were unambiguous. At roughly 130 pressure units right around the maximum service pressure recorded during bench testing the plug threads reached 329~MPa, already above the 310~MPa yield limit of the material. At that same pressure the fuse plate had barely deformed, and the O-ring compression was still sitting around 20\% everywhere along the contour. Even pushing the simulation all the way to 240 units, the fuse still would not open a leak path. \textbf{The fuse was massively oversized and effectively inactive}, leaving the plug as the \emph{de facto} weakest link exactly consistent with the warranty record.

\begin{figure}[H]
    \centering
    \begin{minipage}[t]{0.47\textwidth}
        \centering
        \includegraphics[width=\linewidth]{media/cat/30.png}
    \end{minipage}
    \hfill
    \begin{minipage}[t]{0.47\textwidth}
        \centering
        \includegraphics[width=\linewidth]{media/cat/31.png}
    \end{minipage}
    \caption{Von Mises stress in the fuse assembly at 130 pressure units. The plug threads reach 329~MPa (yield = 310~MPa) while the fuse plate remains elastic, the fuse never triggers before plug failure.}
    \label{fig:cat_fea}
\end{figure}

\subsection{Fuse Re-sizing}

With the fuse proven ineffective, the cheapest corrective action was to keep the same topology but thin down the plate until it actually deformed in the right pressure window. The desired window had clear bounds: seal reliably above the 130-unit maximum service pressure, and leak reliably below the 220-unit cylinder yield pressure a 90-unit band to work inside. (The plug itself is weaker than the cylinder, but its redesign was out of scope for this sizing study; once the plug is redesigned to match the cylinder, the same curves still apply.)

I re-ran the simulation parametrically across four supplier-standard plate thicknesses 3, 3.5, 4 and 6~mm and for each one extracted the minimum and average O-ring compression along the contact contour as pressure ramped. As expected, thinner plates leak earlier and produce wider grey-zone bands. The 6~mm plate essentially replicated the current oversized behaviour; the 3~mm plate leaked too early, barely clearing the 130-unit service ceiling.

\begin{figure}[H]
    \centering
    \includegraphics[width=0.95\textwidth]{media/cat/35.png}
    \caption{O-ring compression vs. pressure for each plate thickness. Grey bands mark the transition between guaranteed sealing ($\min C_\% > 5.7\%$) and certain leakage ($\overline{C_\%} < 0\%$).}
    \label{fig:cat_fusesizing}
\end{figure}

The \textbf{3.5~mm plate} was the clean answer: guaranteed sealing up to around 160 units (comfortably above 130) and guaranteed leakage by around 180 units (comfortably below 220). Part cost stayed negligible since the geometry is still a simple stamped plate. One important limitation worth flagging: this fuse is designed for static or quasi-static overpressure. A sudden pressure spike, for example a hard idler impact  would require a much higher mass-flow evacuation path than the small opening produced by plate deformation. For that failure mode a proper relief valve is the right tool, and that informed the concept-level redesign.

\subsection{Redesign Concepts and Pugh-Matrix Selection}

Three full CAD redesigns were developed in Creo Parametric, each informed by both the warranty analysis and a benchmark of the John~Deere~700L, Komatsu~D51~PX, and CAT~D4/D7/TTL tensioners. The recurring themes across successful competitor designs were consistent: welded (not threaded) cylinder assembly to eliminate the fill-valve leak path, chromed or sleeved piston-cylinder interface to prevent wear-induced seal failure, oil-based lubrication of the sliding contact, a positive mechanical alignment between cylinder and frame, and a proper relief valve as the overpressure safety device rather than a deformable bar.

Concept~1 focused on a perfect-guidance architecture: spherical joint at the piston end, a spacer-and-sleeve retainer providing two line contacts against a precision-machined cylinder, and an interconnected oil chamber lubricating both interfaces. Concept~2 was inspired directly by the John~Deere~700L: a two-piece welded cylinder, a re-sized fuse plate from the sizing study, and a guide hole catching the piston if the chain goes slack. Concept~3 took a different route , it integrates the cylinder body directly into the mobile portion of the TRF, uses a two-level alignment (ball joint at one end, pin-in-hole at the other) that makes the sub-assembly essentially self-aligning under its own weight during installation, and keeps both valves accessible without modifying the TRF hatch door.

The three concepts were scored against eight criteria in a Pugh matrix using the current D5 as the reference: cost, intrusiveness on the existing design, serviceability, operator safety, ease of assembly, behaviour under chain-loose scenarios, fuse robustness, and manufacturability. The matrix was intentionally kept qualitative , weighting the criteria with specific numerical coefficients would have introduced arbitrary bias given that several criteria (assembly ease especially) were the blocking constraints.

\textbf{Concept~3 was selected.} It scored positively on cost, serviceability, manufacturability, intrusiveness and chain-loose guidance; neutral on assembly ease, where its self-aligning geometry directly addresses what had been the single biggest pain point. Concept~2 was blocked by severe assembly difficulty inside the cramped TRF, and Concept~1 offered no clear advantage over the reference. Concept~3 is now positioned for physical prototyping and bench validation the remaining step in the DMAIC Control phase.

\begin{figure}[H]
    \centering
    %------------------------------------------------------------
    % Top block: Concept 1 (left) and Concept 3 (right)
    % Three images stacked; heights fixed so rows align across columns
    %------------------------------------------------------------
    \begin{minipage}[t]{0.48\textwidth}
        \centering
        \includegraphics[width=\linewidth,height=2.4cm,keepaspectratio]{media/cat/37.png}\\[0.5em]
        \includegraphics[width=\linewidth,height=2.4cm,keepaspectratio]{media/cat/38.png}\\[0.5em]
        \includegraphics[width=\linewidth,height=3.8cm,keepaspectratio]{media/cat/39.png}\\[0.7em]
        \textit{Concept 1}
    \end{minipage}\hfill
    \begin{minipage}[t]{0.48\textwidth}
        \centering
        \includegraphics[width=\linewidth,height=2.4cm,keepaspectratio]{media/cat/43.png}\\[0.5em]
        \includegraphics[width=\linewidth,height=2.4cm,keepaspectratio]{media/cat/44.png}\\[0.5em]
        \includegraphics[width=\linewidth,height=3.8cm,keepaspectratio]{media/cat/45.png}\\[0.7em]
        \textit{Concept 3}
    \end{minipage}

    \vspace{1.2em}
    \hrule
    \vspace{1.2em}

    %------------------------------------------------------------
    % Bottom block: Concept 2, three images laid flat
    % Fixed height so they align along a common baseline
    %------------------------------------------------------------
    \begin{minipage}{\textwidth}
        \centering
        \includegraphics[height=3.5cm,keepaspectratio]{media/cat/40.png}\hfill
        \includegraphics[height=3.5cm,keepaspectratio]{media/cat/41.png}\hfill
        \includegraphics[height=3.5cm,keepaspectratio]{media/cat/42.png}\\[0.7em]
        \textit{Concept 2}
    \end{minipage}

    \caption{Comparison of the three design concepts.}
    \label{fig:cat_concepts}
\end{figure}

\begin{tcolorbox}[colback=accentcolor!10,colframe=accentcolor,title=Key Outcomes]
\begin{itemize}[leftmargin=*]
    \item 684 multilingual warranty reports classified automatically at \textbf{96\% accuracy} with a local SBERT pipeline (no data leaves the machine).
    \item FEA with a geometric leak criterion proved the existing fuse never activates, and identified the plug as the real failure point matching field data exactly.
    \item Re-sized fuse plate (3.5~mm) delivers a clean pressure window: seals up to $\sim$160, leaks by $\sim$180, well inside the 130--220 safety band.
    \item Concept~3 selected via Pugh matrix, addressing all four critical failure modes identified upstream and ready for prototyping.
\end{itemize}
\end{tcolorbox}

\newpage

% ============================================
\section{Personal Studies}
% ============================================

\begin{tcolorbox}[colback=gray!10,colframe=primarycolor,title=Note]
These are personal studies aimed at improving my skills in numerical simulation and thermomechanical analysis.
\end{tcolorbox}

\subsection{Study 1: Optimization of Regenerative Cooling in a Combustion Chamber}

\subsubsection{Context}

During operation, rocket engines particularly the combustion chamber are subjected to extremely high temperatures. As temperature increases, the mechanical properties of metals degrade, most notably their yield strength. To address this issue, two methods exist for controlling the temperature of the chamber:

\begin{itemize}[leftmargin=*]
    \item \textbf{Ablative cooling}: involves the gradual erosion of the inner chamber walls, carrying away heat in the process. However, this method is destructive and assumes the engine will not be reused.
    
    \item \textbf{Regenerative cooling}: the most commonly used method. It works by circulating a cryogenic fluid through channels machined into the combustion chamber wall, enabling efficient heat dissipation.
\end{itemize}

\subsubsection{Study Objective}

Evaluate the effect of varying the width $L$ and height $H$ of straight cooling channels on the maximum temperature of the inner combustion chamber wall. The objective is to find an optimum with the maximum inner wall temperature as the objective function and $L$ and $H$ as optimization variables, with:

\begin{itemize}[leftmargin=*]
    \item $L$ between 0.8 and 3.0 mm
    \item $H$ between 1.5 and 5.0 mm
    \item Wall thickness kept constant at 10 mm
    \item Distance between the channel and the inner wall kept constant at 3 mm
\end{itemize}

\begin{figure}[H]
    \centering
    \includegraphics[width=0.9\textwidth]{media/image43.jpg}
    \caption{Cooling channel design (CAD assembly of combustion chamber + nozzle)}
    \label{fig:rocket1}
\end{figure}


\subsubsection{Study Parameters}

\begin{table}[H]
\centering
\begin{tabular}{|l|l|}
\hline
\textbf{Parameter} & \textbf{Value} \\
\hline
Fuel/oxidizer & Methane/liquid oxygen \\
\hline
Inner wall material & CuAgZr \\
\hline
O$_2$/fuel mixture ratio & 3.6 \\
\hline
Chamber pressure (MPa) & 3 \\
\hline
Coolant & N$_2$ (liquid nitrogen) \\
\hline
Coolant mass flow rate (kg/s) & 1.5 \\
\hline
Coolant inlet temperature (K) & 70 \\
\hline
Contraction/expansion ratio & 9/5.5 \\
\hline
Wall thickness (mm) & 10 \\
\hline
Channel-to-wall margin thickness (mm) & 3 \\
\hline
Throat diameter (m) & 0.0274 \\
\hline
\end{tabular}
\caption{Simulation parameters}
\label{tab:rocket_params}
\end{table}

\subsubsection{Execution}

Initially, the simulation was carried out with an inner wall thickness, cooling channel width, and height of 1 mm, 1.8 mm, and 3.5 mm, respectively.

The mesh was generated with inflation layers at the boundary layer of the combustion chamber to capture temperature and velocity gradients, and similarly for the cooling channels. A mesh of the solid region was also generated.

\begin{figure}[H]
    \centering
    \includegraphics[width=1\textwidth]{media/image44.jpg}
    \caption{Fluid + solid zone mesh}
    \label{fig:rocket2}
\end{figure}

After simulation in ANSYS Fluent + ANSYS Thermal, the following results were obtained:

\begin{table}[H]
\centering
\small
\begin{tabular}{|c|c|c|c|c|c|c|}
\hline
\textbf{Axial} & \textbf{Temp.} & \textbf{$\gamma$} & \textbf{Mach} & \textbf{Viscosity} & \textbf{$C_p$} & \textbf{Pr} \\
\textbf{position (mm)} & \textbf{(K)} & & & \textbf{($10^{-4}$ Pa·s)} & \textbf{(J/mol·K)} & \\
\hline
0 & 3448.78 & 1.1251 & 0 & 1.1437 & 2278.4 & 0.5649 \\
\hline
20.5 & 3442.70 & 1.1250 & 0.192 & 1.1422 & 2278.0 & 0.5651 \\
\hline
47.3 & 3287.31 & 1.1213 & 1.000 & 1.1040 & 2267.5 & 0.5716 \\
\hline
79.8 & 2783.91 & 1.1125 & 2.332 & 0.9785 & 2225.3 & 0.5940 \\
\hline
111.64 & 2648.67 & 1.1122 & 2.635 & 0.9444 & 2210.7 & 0.6000 \\
\hline
\end{tabular}
\caption{Thermal data of the gas in the combustion chamber at selected cross-sections}
\label{tab:thermal_data}
\end{table}

\begin{figure}[H]
    \centering
    \includegraphics[width=1\textwidth]{media/image49.png}
    \hfill
    \caption{Temperature distribution in the combustion chamber wall (cross-section taken between two channels)}
    \label{fig:rocket3}
\end{figure}

The temperature peak is reached just before the narrowest cross-section, near the contraction zone of the chamber throat.

After varying parameters $L$ and $H$ in increments of 0.01 mm (between 1.5 and 5.0 mm for height and 0.8 and 3.0 mm for width), the following results were obtained:

\begin{figure}[H]
    \centering
    \includegraphics[width=0.65\textwidth]{media/image50.png}
    \caption{Maximum combustion chamber temperature as a function of L and H}
    \label{fig:rocket4}
\end{figure}

\subsubsection{Conclusion}

\begin{tcolorbox}[colback=accentcolor!10,colframe=accentcolor,title=Optimal Result]
The optimum is reached at $H = 5$ mm and $L = 0.8$ mm, with a minimum temperature of \textbf{586.6 K}.
\end{tcolorbox}

\textbf{Interpretation:} When the channel width is increased, the coolant absorbs a large amount of heat at the inlet of the combustion chamber due to the large exchange surface area, in a region where its effect is not fully needed. By the time it reaches the chamber throat, where temperatures are highest, the coolant has already reached a significant temperature, reducing its ability to absorb heat efficiently.

Reducing the width while increasing the height allows the mass flow rate to be kept constant and delays heat absorption until the throat, where it is most needed.

\newpage

\subsection{Study 2: Failure Prediction of the Starship Tank}

Being a great fan of the \textbf{Starship} rocket project led by SpaceX, I have been closely following its development in detail since 2019. One of the major phases of this development is pressure testing of the LOX (liquid oxygen) and methane tanks.

The goal is to fill the tank with cryogenic liquid nitrogen and increase the internal pressure until rupture, in order to analyze the failure mode. Anticipating this test and having access to a tank drawing published online by an employee including dimensions and sheet thickness. I reconstructed the tank geometry and performed a hydrostatic pressure simulation to attempt to predict the failure zone and the critical pressure using the finite element method.

\begin{figure}[H]
    \centering
    \includegraphics[width=1\textwidth]{media/image51.png}
    \caption{FEA simulation and failure zone prediction}
    \label{fig:starship}
\end{figure}

\begin{tcolorbox}[colback=green!10,colframe=green!50!black,title=Results]
\begin{itemize}[leftmargin=*]
    \item \textbf{Actual tank failure pressure:} 7.6 bar
    \item \textbf{Predicted failure pressure:} 8.1 bar
    \item \textbf{Difference:} Likely due to temperature being higher in the actual tests (the ultimate strength of stainless steel increases as temperature decreases)
    \item \textbf{Failure zone:} Predicted by the red zone (maximum stress). In the actual test, failure occurred exactly in the zone predicted by the finite element method: at the end of the arc and beginning of the linear section.
\end{itemize}
\end{tcolorbox}

Watching the test live and accurately predicting the exact failure zone was particularly striking and validated the simulation approach used.


\newpage

\subsection{Study 3: 6-DOF Rigid-Body Dynamics Simulator, RK4 Integrator}

\subsubsection{Objective}

The goal is simple: obtain a numerically reliable integrator for the six degrees of freedom of a rigid body. No linearisation, no small-angle assumption, no analytical shortcut. Given a force and moment input at each instant, the integrator must propagate the full state of the body forward in time with controlled numerical error and remain stable over long simulation horizons.

\medskip
\noindent\textbf{Full implementation:} \href{https://github.com/floaty-bone/rocket-integrator}{\texttt{github.com/floaty-bone/rocket-integrator}}

\subsubsection{State Representation}

The state of the rigid body is encoded in a 13-element vector:

\[
\mathbf{s} = \bigl[\,\underbrace{x,\,y,\,z}_{\text{position}},\;\underbrace{q_w,\,q_x,\,q_y,\,q_z}_{\text{attitude}},\;\underbrace{v_x,\,v_y,\,v_z}_{\text{velocity}},\;\underbrace{\omega_x,\,\omega_y,\,\omega_z}_{\text{angular velocity}}\,\bigr]
\]

Attitude is represented by a unit quaternion rather than Euler angles. This eliminates gimbal lock entirely and keeps the kinematics well-conditioned at every orientation, including when the body undergoes large rotations.

\subsubsection{Equations of Motion}

The state derivative $\dot{\mathbf{s}}$ is assembled from four coupled equations evaluated at each integration step.

\medskip
\textbf{Position kinematics.}
\[
\dot{\mathbf{r}} = \mathbf{v}
\]

\medskip
\textbf{Quaternion kinematics.} The attitude evolves according to:
\[
\dot{\mathbf{q}} = \frac{1}{2}\,\Omega(\boldsymbol{\omega})\,\mathbf{q}
\]

where $\Omega(\boldsymbol{\omega})$ is the $4\times4$ skew-symmetric matrix constructed from the body-frame angular velocity:

\[
\Omega(\boldsymbol{\omega}) =
\begin{pmatrix}
 0        & -\omega_x & -\omega_y & -\omega_z \\
 \omega_x &  0        &  \omega_z & -\omega_y \\
 \omega_y & -\omega_z &  0        &  \omega_x \\
 \omega_z &  \omega_y & -\omega_x &  0
\end{pmatrix}
\]

\medskip
\textbf{Translational dynamics.} Newton's second law in the inertial frame:
\[
\dot{\mathbf{v}} = \frac{1}{m}\,\mathbf{R}(\mathbf{q})\,\mathbf{F}_{\text{body}} + \mathbf{g}
\]

where $\mathbf{R}(\mathbf{q})$ is the rotation matrix from body frame to inertial frame, derived directly from the quaternion:

\[
\mathbf{R}(\mathbf{q}) =
\begin{pmatrix}
1 - 2(q_y^2+q_z^2)   & 2(q_xq_y - q_wq_z) & 2(q_xq_z + q_wq_y) \\
2(q_xq_y + q_wq_z)   & 1 - 2(q_x^2+q_z^2) & 2(q_yq_z - q_wq_x) \\
2(q_xq_z - q_wq_y)   & 2(q_yq_z + q_wq_x) & 1 - 2(q_x^2+q_y^2)
\end{pmatrix}
\]

\medskip
\textbf{Rotational dynamics.} Euler's rotation equation in the body frame:
\[
\dot{\boldsymbol{\omega}} = \mathbf{I}^{-1}\bigl(\mathbf{M} - \boldsymbol{\omega}\times(\mathbf{I}\,\boldsymbol{\omega})\bigr)
\]

The term $\boldsymbol{\omega}\times(\mathbf{I}\,\boldsymbol{\omega})$ is the gyroscopic term. It accounts for the redistribution of angular momentum between axes whenever the angular velocity vector is not aligned with a principal axis of inertia. It is this term that governs all non-trivial rotational behaviour.

\subsubsection{The RK4 Scheme}

The four equations above are bundled into a single state-derivative function $f(\mathbf{s})$. The classical fourth-order Runge-Kutta method advances the state by one time-step $h$:

\[
\begin{aligned}
\mathbf{k}_1 &= f(\mathbf{s}_n) \\
\mathbf{k}_2 &= f\!\left(\mathbf{s}_n + \tfrac{h}{2}\,\mathbf{k}_1\right) \\
\mathbf{k}_3 &= f\!\left(\mathbf{s}_n + \tfrac{h}{2}\,\mathbf{k}_2\right) \\
\mathbf{k}_4 &= f\!\left(\mathbf{s}_n + h\,\mathbf{k}_3\right) \\[6pt]
\mathbf{s}_{n+1} &= \mathbf{s}_n + \frac{h}{6}\bigl(\mathbf{k}_1 + 2\mathbf{k}_2 + 2\mathbf{k}_3 + \mathbf{k}_4\bigr)
\end{aligned}
\]

The local truncation error is $\mathcal{O}(h^5)$, giving a global error of $\mathcal{O}(h^4)$. At $h = 1\,\text{ms}$, this is more than sufficient for the dynamics of interest.

\textbf{Numerical reliability.} Two measures preserve long-term accuracy. The inverse inertia tensor $\mathbf{I}^{-1}$ is computed once at initialisation and cached, avoiding repeated matrix factorisation at every step. The quaternion is re-normalised after every step to prevent slow numerical drift from corrupting the attitude representation:
\[
\mathbf{q} \leftarrow \frac{\mathbf{q}}{\|\mathbf{q}\|}
\]

\subsubsection{A Fun Observation: the Intermediate Axis Theorem}

Once the integrator was running, the inertia tensor was set to three distinct principal moments, $I_1 = 300$, $I_2 = 100$, $I_3 = 30\,\text{kg}\cdot\text{m}^2$, and the body was given an initial angular velocity nearly aligned with the intermediate axis, with a small perturbation on the other two. The gyroscopic term in Euler's equation then drives the system into an unstable regime: the Dzhanibekov effect. The body undergoes periodic half-turns and re-aligns, indefinitely. The real-time 3-D animation made this directly visible, which is something no closed-form solution could have offered as clearly.

% IMAGE PLACEHOLDER: one frame from the 3-D animation during the Dzhanibekov flip.
% Name the file rk4_dzhanibekov.png, place it in the media/ folder, and uncomment below.
%
% \begin{figure}[H]
%     \centering
%     \includegraphics[width=0.85\textwidth]{media/rk4_dzhanibekov.png}
%     \caption{Real-time 3-D animation: intermediate-axis tumbling (Dzhanibekov effect) reproduced by the RK4 simulator}
%     \label{fig:rk4}
% \end{figure}

\begin{tcolorbox}[colback=accentcolor!10,colframe=accentcolor,title=Key Outcomes]
\begin{itemize}[leftmargin=*]
    \item Full 6-DOF rigid-body propagation with quaternion kinematics, with no gimbal lock.
    \item Fourth-order Runge-Kutta at $h = 1\,\text{ms}$, global error $\mathcal{O}(h^4)$.
    \item Quaternion re-normalisation at every step to prevent long-term attitude drift.
    \item Numerical reproduction of the Dzhanibekov effect: rotation about the intermediate inertia axis is unstable and the simulator captures it exactly.
\end{itemize}
\end{tcolorbox}

\newpage


\subsection{Study 4: Software-in-the-Loop Validation of the LQR Controller}

\subsubsection{Objective}

The LQR controller and the RK4 integrator are two pieces that have been developed and unit-tested in isolation. The controller computes a feedback gain $\mathbf{K}$ from a tangent-space linearisation of the 6-DOF dynamics; the integrator propagates the full non-linear state forward at $1\,\text{ms}$ resolution. Neither result tells us anything about closed-loop behaviour. The Software-in-the-Loop (SIL) simulation closes the loop: at every controller tick, the gain $\mathbf{K}$ is applied to the live plant state, the resulting actuator command is fed back into the RK4 integrator, and the new state is fed back into the controller. The purpose is to validate that the linear LQR law, re-linearised periodically along the trajectory, actually stabilises the full non-linear vehicle from a non-trivial initial attitude to a target setpoint.

\medskip
\noindent\textbf{Full implementation:} \href{https://github.com/floaty-bone/rocket-integrator}{\texttt{github.com/floaty-bone/rocket-integrator}}

\subsubsection{Architecture}

The SIL is built as two decoupled blocks exchanging a strictly defined interface. This boundary is the SIL contract: the controller block knows nothing about the integrator, the plant block knows nothing about the gain matrix, and the bus between them carries only signals that would exist on the real vehicle.

\medskip
\noindent\textbf{Controller block.}
\begin{itemize}[leftmargin=*]
    \item \textbf{Input:} current state $\mathbf{s}\in\mathbb{R}^{13}$, target setpoint $\mathbf{s}^\star\in\mathbb{R}^{13}$.
    \item \textbf{Output:} TVC command $\mathbf{u}_{\text{gimbal}}\in\mathbb{R}^{9}$, encoded as $[\alpha,\,\beta,\,T]$ per engine.
\end{itemize}

\medskip
\noindent\textbf{Plant block.}
\begin{itemize}[leftmargin=*]
    \item \textbf{Input:} current state $\mathbf{s}\in\mathbb{R}^{13}$, TVC command $\mathbf{u}_{\text{gimbal}}\in\mathbb{R}^{9}$.
    \item \textbf{Output:} next state $\mathbf{s}^+\in\mathbb{R}^{13}$, computed by one RK4 step.
\end{itemize}

\medskip
The crucial point is the choice of bus variable. Internally the LQR computes thrust as a Cartesian force vector $[F_x, F_y, F_z]$ per engine, which is the natural form for an affine state-space formulation. The real vehicle, however, commands two gimbal angles and a thrust magnitude. Routing $[\alpha,\,\beta,\,T]$ across the SIL boundary forces the simulation to exercise the same trigonometric mapping the flight software will execute:
\[
\begin{aligned}
F_x &= T\cos\alpha\cos\beta \\
F_y &= T\cos\alpha\sin\beta \\
F_z &= -T\sin\alpha
\end{aligned}
\qquad
\begin{aligned}
T     &= \sqrt{F_x^2+F_y^2+F_z^2} \\
\alpha &= -\arcsin\!\bigl(F_z/T\bigr) \\
\beta  &= \operatorname{atan2}(F_y,\,F_x)
\end{aligned}
\]

Any singularity, saturation, or precision loss in that mapping appears in closed loop, not as a hidden internal quantity.

\subsubsection{Multi-Rate Scheduling}

A real flight stack does not run every block at the same frequency, and the SIL reproduces that. Three independent rates are scheduled inside a single loop:

\begin{table}[H]
\centering
\begin{tabular}{|l|c|l|}
\hline
\textbf{Block} & \textbf{Rate} & \textbf{Justification} \\
\hline
RK4 integration            & 8000 Hz & $h = 0.125\,\text{ms}$, global error $\mathcal{O}(h^4)$ \\
\hline
Controller update          & 5000 Hz & Realistic upper bound for an embedded LQR \\
\hline
LQR re-linearisation       &   30 Hz & JAX auto-diff of $f(\mathbf{s},\mathbf{u})$ and Riccati solve \\
\hline
\end{tabular}
\caption{Scheduled rates inside the SIL loop}
\label{tab:sil_rates}
\end{table}

The relinearisation rate is the load-bearing parameter. LQR is, by definition, a linear law: it is exact only at the operating point $(\mathbf{s}_{\text{op}},\mathbf{u}_{\text{op}})$ where the Jacobians $A=\partial f/\partial\mathbf{s}$ and $B=\partial f/\partial\mathbf{u}$ were taken. The vehicle dynamics are strongly non-linear ($\mathbf{R}(\mathbf{q})$, gyroscopic coupling, $\boldsymbol{\omega}\times(\mathbf{I}\boldsymbol{\omega})$), so the operating point must be refreshed along the trajectory. At each refresh, $A$ and $B$ are recomputed by JAX auto-differentiation of the dynamics function around the current $(\mathbf{s},\mathbf{u})$, and a new gain $\mathbf{K}$ is obtained by solving the continuous-time algebraic Riccati equation:
\[
A^\top P + P A - P B R^{-1} B^\top P + Q = 0,
\qquad
\mathbf{K} = R^{-1} B^\top P
\]

Between refreshes, the controller applies the cached $\mathbf{K}$ at the full 5 kHz rate; the cost stays in the affordable range while the linearisation stays close enough to the trajectory to remain valid.

\subsubsection{Test Scenario}

The vehicle is a 120-tonne stage with three gimballed engines clustered at radius $a = 1.5\,\text{m}$, mounted $l = 18\,\text{m}$ below the centre of mass, with inertia $\operatorname{diag}(1.2\times10^6,\,3.5\times10^7,\,3.5\times10^7)\,\text{kg}\cdot\text{m}^2$. Initial attitude is horizontal $(\theta = -\pi/2)$ at the origin, with zero linear and angular velocity. The setpoint asks for a translation to $(50,\,100,\,60)\,\text{m}$ while holding the same horizontal attitude. The nominal thrust per engine is the static hover share, $mg/3$.

The LQR weights penalise position error heavily ($Q_{\text{pos}} = 10^7$) and linear velocity moderately ($Q_v = 10^5$); attitude and angular rate weights are kept low enough to let the controller redistribute thrust freely. The control cost $R = I_9$ is uniform across the nine actuator channels.

\subsubsection{Results}

The closed loop converges. The vehicle drives all three position channels onto their setpoint with no steady-state offset and no oscillation past the transient. The TVC commands stay inside physically sensible bounds: gimbal angles in the low-degree range, per-engine thrust modulating around the hover share. The final quaternion norm, after 30 seconds at 8 kHz ($2.4 \times 10^5$ RK4 steps), is within $10^{-8}$ of unity, confirming that the per-step re-normalisation of Study 3 holds up over the full SIL horizon.

The most useful diagnostic is the disagreement curve between the controller's internal Cartesian command and the gimbal-encoded command actually delivered to the plant. The two are mathematically inverse on paper; in the SIL they round-trip through floating point, so any divergence flags either a singularity in the $[\alpha,\beta,T]$ chart or a saturation in the controller. Across the test scenario the disagreement stays at machine precision, which is the strongest possible statement that the SIL boundary is exercised cleanly and the controller never asks for a command the real vehicle could not execute.

\begin{tcolorbox}[colback=accentcolor!10,colframe=accentcolor,title=Key Outcomes]
\begin{itemize}[leftmargin=*]
    \item Closed-loop validation of LQR + RK4 across the full SIL boundary at flight-realistic rates (5 kHz control, 8 kHz plant, 30 Hz re-linearisation).
    \item SIL boundary carries $[\alpha,\beta,T]$, the real TVC command, not internal Cartesian forces.
    \item LQR re-linearised on-trajectory via JAX auto-diff + Riccati solve, no hand-derived Jacobians.
    \item Convergence from a horizontal initial attitude to a 3-axis position setpoint, with quaternion norm preserved to $10^{-8}$ after $2.4\times10^5$ integration steps.
\end{itemize}
\end{tcolorbox}

\newpage


\end{document}